\chapter[\emj{🛤️}~Algoritmo de Dijkstra (Camino más Corto)]{\emj{🛤️}~Algoritmo de Dijkstra (Camino más Corto)}

\begin{dsaquees}

El GPS de tu teléfono 📱. Cuando pides la ruta más rápida de tu 
casa al trabajo, el GPS no intenta probar absolutamente todos los
caminos posibles al azar.

En su lugar, siempre expande primero el camino que parece más 
prometedor (el más corto que ha encontrado hasta ahora). Va 
construyendo la ruta óptima paso a paso, actualizando sus cálculos
si encuentra un "atajo" mejor.

Es un algoritmo "Greedy" (Codicioso) que encuentra el camino más corto
desde un nodo origen a todos los demás nodos en un grafo con pesos.

\end{dsaquees}

\begin{dsaotraforma}

Dijkstra resuelve el problema del "Single-Source Shortest Path" en
un grafo ponderado (weighted graph). 

Funciona manteniendo una lista de distancias tentativas y usando un
Min-Heap (Priority Queue) para elegir siempre el nodo con la distancia
mínima actual. 

⚠️ Restricción Vital: NO funciona si el grafo tiene aristas con 
PESOS NEGATIVOS. (Para eso se usa Bellman-Ford).

\end{dsaotraforma}

\begin{dsadiagrama}
\begin{verbatim}
Grafo Ponderado:
       [2]
  (A)-------(B)
   | \       |
(4)|  \(1)   |(3)
   |   \     |
  (C)-------(D)
       (5)

Distancias desde A:
1. Inicio: {A:0, B:∞, C:∞, D:∞}
2. Visitar A: Vecinos B(2), C(4), D(1). 
   Heap elige D(1).
3. Visitar D: Vecinos B(1+3=4), C(1+5=6). 
   ¿4 < dist[B]? No. ¿6 < dist[C]? No.
4. Heap elige B(2).
5. Heap elige C(4).

Final: {A:0, B:2, C:4, D:1}
\end{verbatim}
\end{dsadiagrama}

\begin{lstlisting}[language=C++]
1. Inicializas todas las distancias como Infinito (∞), excepto el
   nodo origen que es 0.
2. Metes (distancia 0, origen) a una Priority Queue (Min-Heap).
3. Mientras la Queue no esté vacía:
   - Sacas el nodo `u` con la menor distancia.
   - Si la distancia sacada es mayor a la que ya conocemos, ignora.
   - Para cada vecino `v` de `u`:
     - Calculas `nueva_dist = dist[u] + peso(u, v)`.
     - SI `nueva_dist < dist[v]`:
       - Actualizas `dist[v] = nueva_dist`.
       - Metes `(nueva_dist, v)` a la Priority Queue.

📝 PSEUDOCÓDIGO
──────────────────────────────────────────────────────────────

función Dijkstra(grafo, origen):
    distancias = array de tamaño V con INFINITO
    distancias[origen] = 0
    pq = crear Min-Heap de pares (distancia, nodo)
    pq.push(0, origen)

    MIENTRAS pq NO vacía:
        (d, u) = pq.pop_min()
        SI d > distancias[u]: continuar
        
        PARA cada vecino v de u con peso w:
            SI distancias[u] + w < distancias[v]:
                distancias[v] = distancias[u] + w
                pq.push(distancias[v], v)

💻 CÓDIGO C++
──────────────────────────────────────────────────────────────

#include <iostream>
#include <vector>
#include <queue>
using namespace std;

const int INF = 1e9;

void dijkstra(int start, vector<vector<pair<int, int>>>& adj) {
    int V = adj.size();
    vector<int> dist(V, INF);
    // Min-heap: priority_queue<T, container, comparison>
    priority_queue<pair<int, int>, vector<pair<int, int>>, greater<pair<int, int>>> pq;

    dist[start] = 0;
    pq.push({0, start});

    while (!pq.empty()) {
        int d = pq.top().first;
        int u = pq.top().second;
        pq.pop();

        if (d > dist[u]) continue;

        for (auto& edge : adj[u]) {
            int v = edge.first;
            int weight = edge.second;

            if (dist[u] + weight < dist[v]) {
                dist[v] = dist[u] + weight;
                pq.push({dist[v], v});
            }
        }
    }

    for(int i = 0; i < V; i++) 
        cout << "Distancia a " << i << ": " << dist[i] << endl;
}

int main() {
    int V = 4;
    vector<vector<pair<int, int>>> adj(V);
    // Definir grafo con pesos...
    dijkstra(0, adj);
    return 0;
}

⏱️ COMPLEJIDAD (Big-O)
──────────────────────────────────────────────────────────────

  Operación        │ Tiempo                │ Razón
  ─────────────────┼───────────────────────┼──────────────────
  Dijkstra         │ O((V + E) log V)      │ Priority Queue
  Espacio          │ O(V + E)              │ Grafo + distancias

* V = Vértices, E = Aristas.
\end{lstlisting}

\begin{dsaentrevistas}

✅ Cuándo usar: "Shortest Path" en grafos con PESOS positivos. Si no
   hay pesos, usa BFS (es más rápido O(V+E)).

💡 Pesos Negativos: Si el entrevistador pregunta "¿Y si hay pesos
   negativos?", responde: "Dijkstra fallaría, usaría Bellman-Ford".

⚠️ C++ Priority Queue: Por defecto es un MAX-heap. Para Dijkstra 
   necesitas un MIN-heap, así que usa `greater<>` o guarda las 
   distancias negativas.

🔑 Variante: "Cheapest Flights Within K Stops" es un problema común 
   de Google que modifica Dijkstra limitando la profundidad.

\end{dsaentrevistas}

\begin{dsaresumen}

• Algoritmo Codicioso (Greedy).
• Usa un Min-Heap para elegir siempre el camino más corto actual.
• Encuentra distancias mínimas desde UN origen a TODOS los demás.
• No admite pesos negativos.
• Tiempo: O(E log V).
• Es la base de los sistemas de navegación modernos.
════════════════════════════════════════════════════════════════

\end{dsaresumen}


